\documentclass[10pt,a4paper]{article}

% Layout: the report must stay within 4 pages (penalty beyond).
\usepackage[margin=2cm]{geometry}
\usepackage[T1]{fontenc}
\usepackage[utf8]{inputenc}
\usepackage[english]{babel}
\usepackage{lmodern}
\usepackage{microtype}

\usepackage{amsmath,amssymb}
\usepackage{graphicx}
\usepackage{booktabs}
\usepackage{subcaption}
\usepackage{xcolor}
\usepackage[numbers,sort&compress]{natbib}
\usepackage[colorlinks=true,linkcolor=blue!50!black,citecolor=blue!50!black,urlcolor=blue!50!black]{hyperref}
\usepackage[capitalize]{cleveref}

\graphicspath{{figures/}}
\setlength{\parskip}{0.3em}

% Notation
\newcommand{\E}{\mathbb{E}}
\newcommand{\Qphi}{Q_{\phi}}
\newcommand{\pith}{\pi_{\theta}}

\title{Overestimation Bias in DDPG and TD3 on LunarLander:\\ the Effect of Layer Normalization}
\author{%
  \begin{tabular}{c}
    Théo Fontaine \qquad Théophile Lahaussois\\[0.4em]
    \small Code: \url{https://github.com/theofontainecs/Overestimation-bias-with-DDPG-and-TD3}
  \end{tabular}}
\date{October 2026}

\begin{document}
\maketitle

\begin{abstract}
This study investigates whether Layer Normalization mitigates or exacerbates the overestimation bias inherent in continuous control algorithms. We evaluate a $2 \times 2$ factorial design (DDPG and TD3, with and without Layer Normalization in the critic) on the LunarLanderContinuous-v3 environment across 5 independent seeds. The overestimation bias is quantified by calculating the mean difference between the critic's Q-value predictions and the true discounted Monte Carlo returns during deterministic policy evaluations. The results show a strong negative correlation between bias and performance; while TD3 successfully bounds overestimation, introducing Layer Normalization into DDPG drastically exacerbates its bias (from $+274.1$ to $+644.8$) and severely degrades final returns.
\end{abstract}

\section{Introduction}
\label{sec:intro}
The combination of function approximation and a deterministic, greedy actor inherently induces an overestimation bias in value estimates \citep{thrun1993issues,fujimoto2018td3}. In baseline continuous control algorithms like DDPG, these positive errors accumulate during temporal difference updates and severely degrade policy performance. To mitigate this structural flaw, TD3 employs a clipped double-Q mechanism, taking the minimum of two critics to strictly bound the target. Concurrently, applying Layer Normalization to the critic is frequently used to stabilize gradients and improve off-policy learning. However, the precise impact of this feature-level normalization on the underlying overestimation bias of these algorithms remains poorly understood. Therefore, this study investigates the following research questions:
\begin{itemize}
  \item[\textbf{Q1}] Does Layer Normalization in the critic reduce the overestimation bias of DDPG and TD3?
  \item[\textbf{Q2}] How does the overestimation bias relate to performance (return)?
\end{itemize}

\section{Background}
\label{sec:background}
\paragraph{DDPG.} \citet{lillicrap2016ddpg}: deterministic actor $\pith$, critic $\Qphi$ trained on
$y = r + \gamma (1-d)\, Q_{\phi'}(s', \pi_{\theta'}(s'))$, actor updated by $\nabla_\theta \E[\Qphi(s,\pith(s))]$.

\paragraph{TD3.} \citet{fujimoto2018td3}: (i) clipped double-Q target
$y = r + \gamma (1-d) \min_{i=1,2} Q_{\phi'_i}(s', \tilde a')$, (ii) target policy smoothing
$\tilde a' = \pi_{\theta'}(s') + \mathrm{clip}(\epsilon, -c, c)$, (iii) delayed actor updates.

\paragraph{Layer Normalization.} \citet{ba2016layernorm} normalise the pre-activations of a layer over its features, per sample, followed by a learned affine map. We insert it after every hidden linear layer of the critic(s), before the ReLU (\texttt{Linear} $\to$ \texttt{LayerNorm} $\to$ \texttt{ReLU}), and nowhere else; the actor is never normalised.

\paragraph{Measuring overestimation.} At each evaluation we run the deterministic policy $\pith$ for 10 episodes and, for every visited pair $(s_t,a_t)$, compare the critic $Q_{\phi_1}$ (the one the actor maximises) with the discounted Monte Carlo return $G_t$, an unbiased estimate of $Q^{\pith}(s_t,a_t)$. We report the mean bias over the $N$ retained pairs, in (discounted) reward units:
\begin{equation}
  \delta = \frac{1}{N}\sum_{t} \big(Q_{\phi_1}(s_t,a_t) - G_t\big), \qquad
  G_t = \sum_{k=0}^{T-t-1} \gamma^k r_{t+k+1},
  \label{eq:bias}
\end{equation}
positive values meaning overestimation. We do not use the normalised bias of \citet{chen2021redq}, which divides by $|\frac1N\sum_t G_t|$: on LunarLander, returns go from negative to positive during learning, and this denominator crosses zero. Episodes truncated by the 1000-step time limit lack the tail of $G_t$, of weight $\gamma^{T-t}$: for them we drop the last $H=\lceil \ln 0.01/\ln\gamma\rceil = 459$ steps, so that the missing tail weighs less than $1\%$. Terminated episodes (landing or crash) are used entirely.

\section{Methodology}
\label{sec:method}
\paragraph{Environment.} \texttt{LunarLanderContinuous-v3} (Gymnasium 1.4): 8-dimensional observations, 2-dimensional actions in $[-1,1]^2$, episodes truncated at 1000 steps; a return of 200 is considered solved.

\paragraph{Implementation.} Our own PyTorch implementation, built on the \texttt{rl-mind} toolkit of the course (environments, replay buffer, collector). DDPG is implemented as TD3 with its three mechanisms switched off (one critic, no target policy smoothing, actor updated at every step); both keep a target actor, as in \citet{fujimoto2018td3}'s DDPG baseline. All other details are shared. Runs were performed on a laptop CPU (Intel i5-10210U), 4 runs in parallel with one thread each, 35 to 42 minutes per run (3\,h\,40 for the 20 runs).

\paragraph{Experimental design.} $2\times 2$ factorial design: \{DDPG, TD3\} $\times$ \{critic without / with LayerNorm\}, 5 seeds per configuration (20 runs), $10^5$ environment steps per run. Every 5000 steps, the deterministic policy is evaluated on 10 episodes (fixed initial-state seeds per run), which give both the evaluation return and the bias of \cref{eq:bias}. Hyper-parameters (\cref{tab:hparams}) are the standard TD3 ones and were not tuned.

\paragraph{Statistics.} Curves show the mean over seeds and a 95\% percentile bootstrap confidence interval (10\,000 resamplings of the seeds). The \emph{final} value of a run is the mean of its last 3 evaluations (90k--100k steps); configurations are compared on these per-seed values with Welch's t-test and the Mann--Whitney U test \citep{colas2019hitchhiker}, at $\alpha = 0.05$. The relation between bias and return is measured by Spearman's rank correlation over runs.

\begin{table}[h]
  \centering
  \small
  \caption{Hyper-parameters, shared by all runs; the last row applies to TD3 only.}
  \label{tab:hparams}
  \begin{tabular}{ll}
    \toprule
    Hidden layers (actor and critics)        & 2 $\times$ 256, ReLU (actor output: tanh) \\
    Optimiser, learning rate (actor, critic) & Adam, $10^{-3}$, $10^{-3}$ \\
    Batch size / replay buffer size          & 256 / $10^5$ (whole history) \\
    Discount $\gamma$ / target update $\tau$ & 0.99 / 0.005 \\
    Exploration                              & 10k uniform random steps, then $\mathcal N(0, 0.1^2)$ \\
    Update-to-data ratio                     & 1 gradient step per environment step \\
    TD3: policy delay / target noise $\sigma$ / clip $c$ & 2 / 0.2 / 0.5 \\
    \bottomrule
  \end{tabular}
\end{table}

\section{Results}
\label{sec:results}

% Figures are produced by 2_analysis.ipynb into report/figures/
\newcommand{\fig}[2]{\IfFileExists{figures/#1.pdf}{\includegraphics[width=#2]{#1.pdf}}{\fbox{\parbox[c][3.5cm][c]{#2}{\centering\ttfamily\detokenize{#1}.pdf}}}}

\begin{figure}[t]
  \centering
  \begin{subfigure}[t]{0.49\linewidth}\centering\fig{returns}{\linewidth}\caption{Evaluation return.}\label{fig:curves}\end{subfigure}\hfill
  \begin{subfigure}[t]{0.49\linewidth}\centering\fig{bias}{\linewidth}\caption{Bias $\delta$ of the critic (\cref{eq:bias}).}\label{fig:bias}\end{subfigure}
  \caption{Learning curves of the four configurations. Lines: mean over 5 seeds; shaded areas: 95\% bootstrap confidence intervals.
Layer Normalization drastically exacerbates the overestimation bias and policy collapse of DDPG, whereas TD3 successfully bounds the bias near zero in both configurations.
  }
  \label{fig:main}
\end{figure}

\subsection{Performance}

As illustrated in \cref{fig:curves}, the baseline DDPG algorithm fails to learn a viable policy, plateauing at a final evaluation return of $-111.0$. Strikingly, the addition of Layer Normalization (LN) to the critic further degrades its performance, pushing the final return down to $-144.1$. In contrast, TD3 demonstrates significantly higher and more stable performance, reaching a final return of $134.6$ and successfully solving the environment in several runs. While the inclusion of LN in TD3 causes a performance drop to $74.2$, it remains vastly superior to the DDPG variants. Overall, the learning curves highlight a clear performance collapse for DDPG, which is only exacerbated by feature-level normalization, whereas TD3 maintains a functional policy throughout training.

\subsection{Overestimation bias}
The performance collapse observed in DDPG can be directly attributed to an underlying overestimation bias, as depicted in \cref{fig:bias} and quantified in \cref{tab:stats}. The baseline DDPG exhibits a substantial final bias of $+274.1$. Critically, the inclusion of LN significantly exacerbates this phenomenon, pushing the bias to an extreme $+644.8$ ($p = 0.016$, Welch and Mann-Whitney U tests). 

Mathematically, this catastrophic increase is rooted in the optimization dynamics governed by the deterministic policy gradient. When LN is applied to the critic's pre-activations, the backward pass scales the gradients by the inverse of the standard deviation ($\frac{\partial \bar{z}}{\partial z} \propto \frac{1}{\sigma}$). This Jacobian preconditioning provides the actor with a perfectly conditioned optimization landscape. Because the critic's estimate inherently contains statistical approximation errors $\epsilon(s,a)$, the perfectly-conditioned actor rapidly converges toward the absolute artificial peaks $\arg\max_a \epsilon(s,a)$. This maximal positive error is then directly bootstrapped into DDPG's unconstrained Bellman target, triggering a fatal positive feedback loop that inflates the bias.

Conversely, TD3 maintains a tightly controlled bias near zero ($-7.8$), and the addition of LN yields only a statistically insignificant variation ($+3.9$, $p = 0.310$). This robustness is achieved through its clipped double-Q mechanism: $y = r + \gamma (1-d) \min_{i=1,2} Q_{\phi'_i}(s', \tilde a')$. Even if LN accelerates the actor's convergence toward local overestimations, the strict $\min$ operator systematically severs the feedback loop by rejecting the most inflated Q-value, completely immunizing the algorithm against the destructive effects of LN.
\begin{table}[h]
  \centering
  \small
  \caption{Final values (mean over 5 seeds of the last 3 evaluations, 95\% bootstrap CI) and pairwise tests between configurations.}
  \label{tab:stats}
  \IfFileExists{figures/stats.tex}{\begin{tabular}{lccc}
\toprule
Configuration & $n$ & Final return [95\% CI] & Final bias [95\% CI] \\
\midrule
DDPG & 5 & -111.0 [-159.6, -63.8] & +274.1 [+203.5, +347.1] \\
DDPG + LN & 5 & -144.1 [-194.0, -78.8] & +644.8 [+466.6, +834.1] \\
TD3 & 5 & 134.6 [37.6, 210.6] & -7.8 [-18.0, +0.6] \\
TD3 + LN & 5 & 74.2 [2.0, 141.1] & +3.9 [-6.5, +15.9] \\
\midrule
Comparison (A vs.\ B) & Metric & $\Delta$ (B $-$ A) & $p$ Welch / MWU \\
\midrule
DDPG vs.\ DDPG + LN & return & -33.1 & 0.468 / 0.548 \\
DDPG vs.\ DDPG + LN & bias & +370.6 & 0.016 / 0.016 \\
TD3 vs.\ TD3 + LN & return & -60.4 & 0.369 / 0.310 \\
TD3 vs.\ TD3 + LN & bias & +11.7 & 0.202 / 0.310 \\
DDPG vs.\ TD3 & return & +245.6 & 0.004 / 0.008 \\
DDPG vs.\ TD3 & bias & -282.0 & 0.002 / 0.008 \\
DDPG + LN vs.\ TD3 + LN & return & +218.3 & 0.003 / 0.016 \\
DDPG + LN vs.\ TD3 + LN & bias & -640.9 & 0.003 / 0.008 \\
\bottomrule
\end{tabular}
}{}
\end{table}

\subsection{Bias vs.\ performance}
\begin{figure}[h]
  \centering
  \fig{bias_vs_return}{0.55\linewidth}
  \caption{Final bias vs.\ final return, one point per run (20 runs).}
  \label{fig:scatter}
\end{figure}
\cref{fig:scatter} illustrates the direct relationship between the final bias and the final evaluation return across all 20 independent runs. We observe a strong, highly significant negative rank correlation (Spearman $\rho = -0.87$, $p < 0.001$). The scatter plot clearly segregates the algorithms into two distinct clusters: the TD3 runs (with and without LN) are tightly grouped in the top-left quadrant, characterized by high returns and near-zero bias. In stark contrast, the DDPG variants are scattered across the bottom-right quadrant, suffering from severe overestimation and poor performance. This structural segregation confirms empirically that overestimation is not merely a numerical artifact, but the primary driver of policy failure in unconstrained architectures.
\section{Discussion}
\label{sec:discussion}
These statistical findings provide clear answers to our initial research questions. Regarding \textbf{Q1}, Layer Normalization does not reduce overestimation bias; rather, it drastically amplifies it in unconstrained architectures like DDPG, while TD3's clipped double-Q mechanism effectively blocks this amplification. Regarding \textbf{Q2}, the strong negative correlation ($\rho = -0.87$) confirms that overestimation bias is not a benign artifact, but a primary driver of policy collapse. 

However, these conclusions are subject to methodological limitations. The empirical evidence relies on a single continuous control task (LunarLanderContinuous) and a restricted sample size (5 seeds per configuration), which limits generalizability to higher-dimensional environments. Furthermore, the reliance on Monte Carlo returns ($G_t$) to estimate the true Q-value introduces high variance into the bias measurement $\delta$ (as defined in \cref{eq:bias}), which may partially noise the evaluation despite the averaging over $N$ pairs.

\section{Conclusion}
\label{sec:conclusion}
This study demonstrates that while Layer Normalization is a powerful architectural tool for gradient conditioning, it cannot substitute for fundamental algorithmic safeguards. By providing a perfectly conditioned optimization landscape, LN inadvertently accelerates the exploitation of approximation errors in DDPG, leading to catastrophic overestimation. Consequently, our findings underscore the absolute necessity of structural pessimism—such as TD3's minimum operator—to safely deploy feature-level normalization in continuous deep reinforcement learning.

\bibliographystyle{plainnat}
\bibliography{refs}

\end{document}
